\documentclass[12pt,a4paper]{article}
\usepackage[utf8]{vietnam}
\usepackage{amsmath,amssymb}
\usepackage{tabularx}
\usepackage{tikz}
\usetikzlibrary{calc,angles,quotes,arrows.meta,patterns}
\usepackage{xcolor}
\usepackage{enumitem}
\usepackage{fancyhdr}
\usepackage{graphicx}
\usepackage[top=1.0cm, bottom=1.0cm, left=1.2cm, right=1.2cm]{geometry}
\usepackage[dethi]{ex_test}
\SAOPTN{widthSA=3.2,heightSA=0.62}

\setlength{\headheight}{15pt}
\setlength{\parskip}{0pt}

\pagestyle{fancy}
\fancyhf{}
\rhead{\textit{Môn Toán 8 -- Đề số 03}}
\lhead{\textit{Bộ đề ôn tập Giữa học kì I}}
\cfoot{\thepage}

\begin{document}

% --- KHUNG TIÊU ĐỀ ĐỀ THI ---
\noindent
\begin{tabularx}{\linewidth}{X r}
\textbf{TRƯỜNG THCS: \dotfill} & \textbf{ĐỀ KIỂM TRA ĐỊNH KỲ GIỮA HỌC KỲ I} \\
\textbf{Họ và tên: \dotfill} & \textbf{NĂM HỌC 2026 -- 2027} \\
\textbf{Lớp: \dotfill \quad SBD: \dotfill} & \textbf{Môn: TOÁN 8} \\
\hline
\end{tabularx}
\vspace{2mm}

\begin{center}
  \textbf{\large ĐỀ ÔN TẬP GIỮA HỌC KỲ I -- ĐỀ SỐ 03}\\
  \textit{Thời gian làm bài: 90 phút (Không kể thời gian phát đề)}
\end{center}
\vspace{1mm}

% --- PHẦN I ---
\noindent\textbf{\large PHẦN I. CÂU TRẮC NGHIỆM NHIỀU PHƯƠNG ÁN LỰA CHỌN (3,0 điểm)}

\textit{Thí sinh trả lời từ câu 1 đến câu 12. Mỗi câu hỏi thí sinh chỉ chọn một phương án duy nhất.}

\vspace{1mm}
\begin{ex}
Biểu thức nào sau đây là đơn thức?
\choice
{$x^3(1 + x)$}
{\True $3xy^2z^5$}
{$(xy + z)t$}
{$\dfrac{2}{x} + y$}
\end{ex}

\begin{ex}
Cặp đơn thức nào sau đây là đồng dạng với nhau?
\choice
{$\dfrac{1}{2}x^2y$ và $2xy^2$}
{$3xy$ và $-5x$}
{\True $x^2y^4$ và $2x^2y^4$}
{$-3xy$ và $-3xy^2$}
\end{ex}

\begin{ex}
Thực hiện phép chia đơn thức $5x^2y^2z$ cho $5xyz$ ta được kết quả là:
\choice
{$5xy$}
{\True $xy$}
{$yz$}
{$25x^3y^3z^2$}
\end{ex}

\begin{ex}
Hằng đẳng thức $A^2 - B^2 = (A - B)(A + B)$ có tên là:
\choice
{Bình phương của một tổng}
{Bình phương của một hiệu}
{Tổng hai bình phương}
{\True Hiệu hai bình phương}
\end{ex}

\begin{ex}
Tứ giác nào dưới đây là hình thoi?
\choice
{\True Tứ giác có 4 cạnh bằng nhau}
{Tứ giác có hai đường chéo vuông góc}
{Hình bình hành có hai đường chéo bằng nhau}
{Tứ giác có hai đường chéo cắt nhau tại trung điểm mỗi đường}
\end{ex}

\begin{ex}
Hai đường chéo của hình chữ nhật có tính chất nào sau đây?
\choice
{Chúng vuông góc với nhau}
{Chúng bằng nhau}
{Chúng cắt nhau tại trung điểm của mỗi đường}
{\True Chúng bằng nhau và cắt nhau tại trung điểm của mỗi đường}
\end{ex}

\begin{ex}
Hình bình hành $ABCD$ có $\widehat{A} = 65^\circ$. Số đo góc $\widehat{C}$ là:
\choice
{$125^\circ$}
{$95^\circ$}
{$115^\circ$}
{\True $65^\circ$}
\end{ex}

\begin{ex}
Cho tam giác $ABC$ vuông tại $A$ có $M$ là trung điểm của $BC$, biết $BC = 10\text{ cm}$. Khi đó độ dài đoạn thẳng $AM$ bằng:
\choice
{\True $5\text{ cm}$}
{$10\text{ cm}$}
{$15\text{ cm}$}
{$20\text{ cm}$}
\end{ex}

\begin{ex}
Cho các tam giác có độ dài ba cạnh như sau. Tam giác nào là tam giác vuông?
\choice
{$15\text{ cm}; 8\text{ cm}; 10\text{ cm}$}
{$6\text{ dm}; 8\text{ dm}; 9\text{ dm}$}
{\True $5\text{ cm}; 4\text{ cm}; 3\text{ cm}$}
{$2\text{ cm}; 3\text{ cm}; 4\text{ cm}$}
\end{ex}

\begin{ex}
Cho hình chữ nhật $ABCD$. Đoạn thẳng $AC$ bằng đoạn thẳng nào sau đây?
\choice
{$AB$}
{$BC$}
{$AD$}
{\True $BD$}
\end{ex}

\begin{ex}
Cho tứ giác $ABCD$, trong đó $\widehat{A} + \widehat{B} = 140^\circ$. Tổng $\widehat{C} + \widehat{D}$ bằng:
\choice
{\True $220^\circ$}
{$200^\circ$}
{$160^\circ$}
{$130^\circ$}
\end{ex}

\begin{ex}
Thể tích của một hình chóp tam giác đều có diện tích đáy bằng $15\text{ cm}^2$ và chiều cao bằng $6\text{ cm}$ là:
\choice
{$90\text{ cm}^3$}
{\True $30\text{ cm}^3$}
{$45\text{ cm}^3$}
{$15\text{ cm}^3$}
\end{ex}


\vspace{1mm}
\noindent\textbf{BẢNG TRẢ LỜI TRẮC NGHIỆM PHẦN I:}
\begin{center}
\begin{tabular}{|c|c|c|c|c|c|c|c|c|c|c|c|c|}
\hline
\textbf{Câu} & \textbf{1} & \textbf{2} & \textbf{3} & \textbf{4} & \textbf{5} & \textbf{6} & \textbf{7} & \textbf{8} & \textbf{9} & \textbf{10} & \textbf{11} & \textbf{12} \\
\hline
\textbf{Chọn} & \rule{0pt}{0.65cm} & & & & & & & & & & & \\
\hline
\end{tabular}
\end{center}
\vspace{2mm}

\setcounter{ex}{0}
\noindent\textbf{\large PHẦN II. CÂU TRẮC NGHIỆM ĐÚNG SAI (2,0 điểm)}

\textit{Thí sinh trả lời câu 1 và câu 2. Trong mỗi ý a), b), c), d) ở mỗi câu, thí sinh chọn đúng hoặc sai.}

\vspace{1mm}
\begin{ex}
Cho hai biểu thức $A = (2x - 1)(x + 3) - 2x(x + 2)$ và $B = (6x^3y^2 - 9x^2y^3) : 3x^2y$. Xét tính đúng sai của các mệnh đề sau:
\choiceTF
{\True Khai triển tích $(2x - 1)(x + 3)$ ta được $2x^2 + 5x - 3$}
{\True Kết quả rút gọn của biểu thức $A$ là $x - 3$}
{\True Kết quả của phép chia biểu thức $B$ là $2xy - 3y^2$}
{Tại $x = 2$ và $y = 1$, giá trị của biểu thức $A$ gấp 2 lần giá trị của biểu thức $B$}
\end{ex}

\begin{ex}
Cho hai đa thức $P = x^2 - 6x + 9 - y^2$ và $Q = 8x^3 - 27$. Xét tính đúng sai của các mệnh đề sau:
\choiceTF
{\True Biến đổi $x^2 - 6x + 9$ thành bình phương của một hiệu ta được $(x - 3)^2$}
{\True Phân tích đa thức $P$ thành nhân tử ta được $(x - y - 3)(x + y - 3)$}
{\True Phân tích đa thức $Q$ thành nhân tử theo hằng đẳng thức hiệu hai lập phương ta được $(2x - 3)(4x^2 + 6x + 9)$}
{Tại $x = 5$ và $y = 2$, giá trị của đa thức $P$ bằng 8}
\end{ex}


\noindent\textbf{BẢNG TRẢ LỜI PHẦN II:}
\begin{center}
\begin{tabular}{|c|c|c|c|c||c|c|c|c|c|}
\hline
\textbf{Câu 1} & \textbf{a)} & \textbf{b)} & \textbf{c)} & \textbf{d)} & \textbf{Câu 2} & \textbf{a)} & \textbf{b)} & \textbf{c)} & \textbf{d)} \\
\hline
\textbf{Đ/S} & \rule{0pt}{0.6cm} & & & & \textbf{Đ/S} & & & & \\
\hline
\end{tabular}
\end{center}
\vspace{2mm}

\setcounter{ex}{0}
\noindent\textbf{\large PHẦN III. CÂU HỎI TRẢ LỜI NGẮN (2,0 điểm)}

\textit{Thí sinh trả lời câu 1 và câu 2. Điền kết quả vào ô trống tương ứng.}

\vspace{1mm}
\begin{ex}
Rút gọn biểu thức: $A = (2x - 1)(x + 3) - 2x(x + 2)$. \shortans{$x - 3$}
\end{ex}

\begin{ex}
Phân tích đa thức sau thành nhân tử: $x^2 - y^2 + 2x + 1$. \shortans{$(x - y + 1)(x + y + 1)$}
\end{ex}


\setcounter{ex}{0}
\noindent\textbf{\large PHẦN IV. TỰ LUẬN (3,0 điểm)}

\textit{Thí sinh trình bày chi tiết lời giải vào giấy làm bài.}

\vspace{1mm}
\begin{ex} (2,0 điểm)
Thực hiện phép tính:
\begin{enumerate}[label=\textbf{\alph*)}, leftmargin=0.8cm, itemsep=2pt]
\item (0,5 điểm) Thu gọn và tính giá trị biểu thức $A = 5x^2y - 2xy + 3x^2y - xy$ tại $x = -1, y = 2$.
\item (0,75 điểm) Thực hiện phép nhân: $B = 3xy(2x^2 - 5y^3)$.
\item (0,75 điểm) Thực hiện phép chia: $C = (6x^3y^2 - 9x^2y^3) : 3x^2y^2$.
\end{enumerate}
\end{ex}
\vspace{0.5mm}

\begin{ex} (1,5 điểm)
Cho hai đa thức: $P = 2x^2y + 5x^3 - 2xy^2 + 6$ và $Q = x^3 + xy^2 - x^2y - 2$.
\begin{enumerate}[label=\textbf{\alph*)}, leftmargin=0.8cm, itemsep=2pt]
\item (0,75 điểm) Tính tổng của hai đa thức $P$ và $Q$.
\item (0,75 điểm) Tìm đa thức $N$ sao cho $Q = P + N$.
\end{enumerate}
\end{ex}
\vspace{0.5mm}

\begin{ex} (1,0 điểm)
Khu vườn của nhà bác Hoa có dạng hình vuông. Bác Hoa muốn dành một mảnh đất hình chữ nhật ở góc khu vườn làm nhà để dụng cụ làm vườn (hình vẽ).
\begin{center}
\begin{tikzpicture}[scale=0.8, baseline=(current bounding box.center)]
  \draw[thick, fill=green!10] (0,0) rectangle (4,4);
  \draw[thick, fill=orange!25] (0,0) rectangle (2.2,1.3);
  \node at (1.1,0.65) {\small Nhà kho};
  \node at (2.5,2.6) {\small Khu vườn};
  \draw[<->] (0,-0.3) -- (4,-0.3) node[midway, below] {\small $x$};
  \draw[<->] (-0.3,0) -- (-0.3,4) node[midway, left] {\small $x$};
  \draw[<->] (2.2,-0.3) -- (4,-0.3) node[midway, below] {\small $15\text{ m}$};
  \draw[<->] (-0.3,1.3) -- (-0.3,4) node[midway, left] {\small $25\text{ m}$};
\end{tikzpicture}
\end{center}
\begin{enumerate}[label=\textbf{\alph*)}, leftmargin=0.8cm, itemsep=2pt]
\item (0,5 điểm) Viết đa thức biểu thị chu vi của mảnh đất làm nhà.
\item (0,5 điểm) Biết chu vi của mảnh đất làm nhà bằng $40\text{ m}$. Tính diện tích khu vườn hình vuông ban đầu.
\end{enumerate}
\end{ex}
\vspace{0.5mm}

\begin{ex} (0,5 điểm)
Cho các số thực $x, y$ thỏa mãn $x + y = 1$. Tìm giá trị nhỏ nhất của biểu thức $A = x^3 + y^3 + xy$.
\begin{enumerate}[label=\textbf{\alph*)}, leftmargin=0.8cm, itemsep=2pt]
\item (0,5 điểm) Tìm GTNN của $A$.
\end{enumerate}
\end{ex}
\vspace{0.5mm}


\begin{center}
\rule{4cm}{0.4pt}\\
\textit{--- HẾT ---}
\end{center}
\end{document}
